\documentclass[10pt]{article}
\usepackage[UTF8]{ctex}
\usepackage{amsmath, amssymb}
\usepackage{geometry}
\geometry{a4paper, margin=1.5cm}
\usepackage{enumitem}

\title{计数原理入门}
\author{}
\date{}

\begin{document}
\maketitle

\section{导言：为什么要学计数？}

计数原理研究的是：\textbf{在给定条件下，所有可能的不同结果一共有多少种}。
它不关心具体哪个结果会发生，只关心可能性的总数。这种“不重不漏”地计算总数的方法，是概率、统计、组合数学的基础。

你将学到两个最基本的计数规则（加法原理、乘法原理），以及它们最重要的两个应用模型：\textbf{排列}（考虑顺序）和\textbf{组合}（不考虑顺序）。

\section{两个基本计数原理}

\subsection{分类加法原理}

\textbf{表述}
完成一件事有 \(k\) 类不同的方案。
第 1 类有 \(m_1\) 种方法，第 2 类有 \(m_2\) 种方法，……，第 \(k\) 类有 \(m_k\) 种方法。
并且任意两类方案之间没有重叠（即一个结果只属于其中一类）。
那么完成这件事共有
\[
N = m_1 + m_2 + \dots + m_k
\]
种方法。

\textbf{关键词}：分类、加法、“或”。

\textbf{逻辑本质}
如果完成目标的各个途径\textbf{相互独立}（选了这个就不能选那个），那么总数为各类方法数之和。
这源于集合的\textbf{不相交并集}的基数公式：若 \(A_1, A_2, \dots, A_k\) 两两不交，则 \(|A_1 \cup \cdots \cup A_k| = |A_1| + \cdots + |A_k|\)。

\textbf{例 1}：从某城市到机场：可以乘地铁（3条线路），或乘机场大巴（2条线路），或乘出租车（无数种，但视为“1类方法”——只要知道有出租车这一选项即可）。
若只考虑地铁和机场大巴，那么不同走法总数为 \(3 + 2 = 5\)。
这里“乘地铁”与“乘机场大巴”是两类互斥的选择。

\subsection{分步乘法原理}

\textbf{表述}
完成一件事需要依次经过 \(k\) 个步骤。
第 1 步有 \(m_1\) 种做法，第 2 步有 \(m_2\) 种做法，……，第 \(k\) 步有 \(m_k\) 种做法。
并且\textbf{各步的做法相互独立}（即第 1 步的每一种选择，都不影响第 2 步的可选方法数，依此类推）。
那么完成这件事共有
\[
N = m_1 \times m_2 \times \dots \times m_k
\]
种方法。

\textbf{关键词}：分步、乘法、“且”。

\textbf{逻辑本质}
完成整个过程等价于依次选择第一步的结果、第二步的结果……每一步的选择构成一个有序 \(k\) 元组。
总数为各步方法数的乘积，这来自\textbf{笛卡尔积}的基数公式：\(|A \times B| = |A| \cdot |B|\)。

\textbf{例 2}：设置一个两位数字密码：第一位从 1~9 中选（9种），第二位从 0~9 中选（10种），且选择独立。
那么不同密码总数为 \(9 \times 10 = 90\)。
这里“先选第一位，再选第二位”是两步，缺一不可。

\subsection{两个原理的对比与联系}

\begin{tabular}{|l|l|l|}
\hline
比较维度 & 加法原理 & 乘法原理 \\
\hline
完成方式 & 多类方案中\textbf{任选一类}即可完成 & 必须\textbf{依次完成所有步骤}才算完成 \\
\hline
逻辑关系 & 或 & 且 \\
\hline
运算 & 加 & 乘 \\
\hline
集合论基础 & 不相交并集 & 笛卡尔积 \\
\hline
\end{tabular}

\textbf{常见混淆}：
如果一个问题的描述中出现“或者”，往往要用加法；出现“并且”“然后”，往往要用乘法。
但更可靠的方法是：问自己“能否将整个过程拆成相互独立的先后阶段？”——能，则乘法；不能，则看是否分成互斥的类别。

\section{排列与组合：两个重要的计数模型}

当问题涉及\textbf{从有限个不同元素中选出一部分}时，根据选出的元素\textbf{是否考虑顺序}，我们得到两个不同的概念：排列和组合。

\subsection{排列（有顺序的选取）}

\textbf{定义}
从 \(n\) 个\textbf{不同}的元素中，取出 \(m\) 个（\(1 \le m \le n\)），并按照一定的顺序排成一列，称为从 \(n\) 个元素中取出 \(m\) 个元素的\textbf{一个排列}。
所有不同排列的个数称为\textbf{排列数}，记作 \(A_n^m\)（或 \(P_n^m\)）。

\textbf{公式推导（使用乘法原理）}
排第一个位置：有 \(n\) 种选择；
排第二个位置：从剩下的 \(n-1\) 个中选，有 \(n-1\) 种；
……
排第 \(m\) 个位置：从剩下的 \(n - (m-1) = n-m+1\) 个中选，有 \(n-m+1\) 种。
因此
\[
A_n^m = n \times (n-1) \times \dots \times (n-m+1).
\]
用阶乘表示：
\[
A_n^m = \frac{n!}{(n-m)!}, \quad \text{其中 } n! = n \times (n-1) \times \cdots \times 2 \times 1,\ 0! = 1.
\]

\textbf{关键特征}：顺序不同视为不同结果。
例如：从 \{a, b, c\} 中选 2 个排列，ab 与 ba 是两个不同的排列。

\subsection{组合（无顺序的选取）}

\textbf{定义}
从 \(n\) 个不同元素中取出 \(m\) 个（\(0 \le m \le n\)），并成一组（不考虑组内顺序），称为从 \(n\) 个元素中取出 \(m\) 个元素的\textbf{一个组合}。
所有不同组合的个数称为\textbf{组合数}，记作 \(C_n^m\)（或 \(\binom{n}{m}\)）。

\textbf{公式推导（利用排列与组合的关系）}
先从 \(n\) 个元素中选 \(m\) 个进行排列，有 \(A_n^m\) 种。
但对于同一个组合（即同一组元素），将它们进行全排列会得到 \(m!\) 个不同的排列。
因此：
\[
A_n^m = C_n^m \times m! \quad \Rightarrow \quad C_n^m = \frac{A_n^m}{m!} = \frac{n!}{m!\,(n-m)!}.
\]

\textbf{关键特征}：顺序无关。
例如：从 \{a, b, c\} 中选 2 个组合，\{a, b\} 和 \{b, a\} 是同一个组合。

\subsection{排列与组合的对比}

\begin{tabular}{|l|l|l|}
\hline
对比项 & 排列 \(A_n^m\) & 组合 \(C_n^m\) \\
\hline
顺序是否重要 & 是 & 否 \\
\hline
公式 & \(\frac{n!}{(n-m)!}\) & \(\frac{n!}{m!(n-m)!}\) \\
\hline
数值关系 & \(A_n^m = C_n^m \times m!\) & \(C_n^m = A_n^m / m!\) \\
\hline
典型表述 & “排成一排”“安排职位”“顺序不同即不同” & “选出一组”“抽取样品”“只关心哪些元素” \\
\hline
\end{tabular}

\textbf{注意}：当 \(m=0\) 时，规定 \(A_n^0 = 1\)（一种排列：什么都不排），\(C_n^0 = 1\)（一种组合：什么都不选）。

\section{符号体系与基本运算工具}

\subsection{阶乘 \(n!\)}

\begin{itemize}
\item 定义：\(n! = n \times (n-1) \times \cdots \times 2 \times 1\)，特别地 \(0! = 1\)。
\item 性质：\((n+1)! = (n+1) \times n!\)。
\item 作用：简洁表达排列数和组合数。
\end{itemize}

\textbf{示例}
\(5! = 120\)，\(4! = 24\)，\(\frac{5!}{3!} = 5 \times 4 = 20 = A_5^2\)。

\subsection{排列数符号 \(A_n^m\)}

\begin{itemize}
\item 读作“\(n\) 选 \(m\) 的排列数”。
\item 注意：\(A_n^m\) 中 \(n \ge m \ge 0\)，且 \(n,m\) 均为整数。
\item 性质：\(A_n^n = n!\)（全排列）。
\end{itemize}

\subsection{组合数符号 \(C_n^m\)}

\begin{itemize}
\item 读作“\(n\) 选 \(m\) 的组合数”。
\item 常用性质（要求了解，但不需在此深入证明）：
  \begin{itemize}
  \item 对称性：\(C_n^m = C_n^{n-m}\)。
  \item 递推公式：\(C_n^m = C_{n-1}^{m-1} + C_{n-1}^{m}\)（杨辉三角的构造依据）。
  \end{itemize}
\end{itemize}

\subsection{求和符号 \(\sum\)}

\begin{itemize}
\item 形式：\(\sum_{i=1}^{k} a_i = a_1 + a_2 + \cdots + a_k\)。
\item 在计数中用于表达加法原理的总数。
\end{itemize}

\section{蕴含的数学思想（重点）}

\subsection{分类讨论思想}

当问题不能用一个统一的乘法过程解决时，需要根据某种属性将结果分成\textbf{互不相交}的类别，分别计算每类的方法数，再求和。
关键要求：\textbf{不重不漏}。
这也是解决复杂计数问题的基本策略之一。

\subsection{有序与无序的区分}

\begin{itemize}
\item \textbf{有序}：结果由“选哪些”和“以什么顺序”共同决定 → 排列模型。
\item \textbf{无序}：结果仅由“选哪些”决定 → 组合模型。
\end{itemize}
区分二者的能力，是正确使用排列或组合公式的前提。

\subsection{化归思想（转化为基本模型）}

许多表面复杂的问题，可以通过适当分解，归为若干个加法原理或乘法原理的简单组合。
例如：先分类，再对每一类分步；或者先分步，再对某一步进行分类。
这就是“先选后排”“特殊位置优先”等技巧的思想根源。

\subsection{符号化与公式化}

用 \(A_n^m, C_n^m\) 等符号代替冗长的连乘表达式，使思考聚焦于\textbf{结构}而非具体数值。
公式是规律的浓缩，学会识别问题属于哪个模型，直接代入公式，是提高效率的关键。

\section*{结语：接下来学什么？}

你已经掌握了计数原理的\textbf{核心概念}：
\begin{itemize}
\item 何时用加法、何时用乘法
\item 排列与组合的区别与联系
\item 常用符号与基本公式
\item 重要的数学思想（分类、有序/无序、化归）
\end{itemize}

这些知识足以支撑你阅读绝大多数计数问题的题目表述。
下一步将是\textbf{练习}：通过具体问题训练自己识别模型、正确使用公式、处理约束条件（如“不相邻”“至少一个”等）。
但在此之前，请反复理解本读物中的定义和思想——它们是解题的“语法规则”，比刷题更重要。

\begin{quote}
\textbf{说明：} 这份读物不追求趣味性，而追求概念的准确性和逻辑的清晰性。建议你边读边在草稿纸上写下关键公式，并自己举例验证每个定义。
\end{quote}



\end{document}